\documentclass{ohm_project_description}
\usepackage[english]{babel}

\projecttitle{M-APR: Emergent Systems in Production}
\projectauthor{TTZ Nürnberger Land / TH Nürnberg}
\projectdate{2026}

\begin{document}

\maketitle
\thispagestyle{fancy}

Modern production environments increasingly rely on heterogeneous robot fleets — autonomous mobile robots (AMR), collaborative robot arms and stationary processing stations — that jointly handle manufacturing and logistics tasks. Today, coordination is predominantly centralised. However, centralised coordination is vulnerable to disruptions: if the central node fails, the entire process chain comes to a halt.

This M-APR project (3 semesters, starting WiSe 2026/2027) investigates whether and under what conditions decentralised, emergent coordination strategies outperform centralised control in production scenarios. The focus is on the systematic comparison of both approaches under controlled variations of adversarial conditions (communication failures, machine failures, dynamic order changes). In the third semester, selected scenarios are transferred to real hardware at the TTZ Nürnberger Land.

\section*{Work Packages (3 Semesters)}
\begin{itemize}[leftmargin=0.5cm]
    \setlength\itemsep{.1em}
    \item \textbf{Semester 1:} Literature review, ROS\,2 simulation environment, centralised baseline and first decentralised approach
    \item \textbf{Semester 2:} Multi-Agent Reinforcement Learning (MARL), systematic comparison under adversarial conditions, statistical evaluation
    \item \textbf{Semester 3:} Transfer to real hardware at TTZ, sim-to-real evaluation, master's thesis and conference publication
\end{itemize}

\section*{Requirements}
\begin{itemize}[leftmargin=0.5cm]
    \setlength\itemsep{.1em}
    \item Bachelor's degree in electrical engineering, computer science, mechatronics, robotics or equivalent
    \item Programming in C++ and/or Python
    \item Basic knowledge of machine learning / reinforcement learning
    \item Experience with ROS\,2 and simulation (Gazebo or similar) is an advantage
\end{itemize}

\vspace{0.5cm}
This topic is designed as an \textbf{M-APR (3 semesters)}, but can also be completed as a project or master's thesis subject to agreement.

\vfill
\textcolor{ohm_red}{\rule{\linewidth}{0.4mm}}
\textbf{\textcolor{ohm_red}{TTZ Nürnberger Land / Mobile Robotics Lab}} \\
\begin{tabular}{@{}ll}
\textbf{Supervisor:}    & Prof. Dr. Christian Pfitzner \\
\textbf{E-Mail:}        & \href{mailto:christian.pfitzner@th-nuernberg.de}{christian.pfitzner@th-nuernberg.de} \\
\textbf{Co-Supervisor:} & Waldemar Haag \\
\textbf{E-Mail:}        & \href{mailto:waldemar.haag@th-nuernberg.de}{waldemar.haag@th-nuernberg.de} \\
\end{tabular}

\end{document}
